\documentclass[aspectratio=169]{beamer}
\usetheme{Madrid}
\usepackage[utf8]{inputenc}
\usepackage[T1]{fontenc}
\usepackage{tikz}
\usepackage{amssymb}
\usepackage{graphicx}
\setbeamertemplate{navigation symbols}{}
\setbeamercolor{title}{fg=white,bg=magenta!70!red}
\setbeamercolor{frametitle}{fg=white,bg=magenta!70!red}
\setbeamercolor{block title}{bg=magenta!70!red!20,fg=black}
\setbeamercolor{block body}{bg=magenta!70!red!8,fg=black}
\setbeamercolor{item}{fg=magenta!70!red}
\setbeamercolor{itemize item}{fg=magenta!70!red}
\setbeamertemplate{itemize item}[triangle]

\title{Foundations 7 Capstone — Flight Computer Online}
\subtitle{A loop AND a decision, working together on real sensor data}
\author{ByteForge --- Grade 9 --- The Deep Diver}
\date{}

\begin{document}

\begin{frame}
    \titlepage
\end{frame}

\begin{frame}{Today's Missions}
    \Large
    \begin{itemize}
        \item BOSS: Flight Computer Online
    \end{itemize}
\end{frame}

\begin{frame}{Every real algorithm combines simple pieces}
    \large
    \begin{itemize}
        \item Loop through every sensor reading (Foundations 7)
        \item Make a decision per item --- is THIS reading above the threshold?
        \item Count how many readings cross it --- a genuinely common real data-processing shape
    \end{itemize}
\end{frame}

\begin{frame}{Real-World Examples}
    \large
    \begin{itemize}
        \item A weather station counting "how many days this month were over 30 degrees" is exactly this pattern: loop + threshold decision + counter.
    \end{itemize}
\end{frame}


\begin{frame}{Let Them Discover}
    \large
    Given readings 80, 90, 70, 85, 95 and a threshold of "above 80", work out by hand how many actually count --- then check your answer against the mission.
\end{frame}

\begin{frame}{Your Turn}
    \Large
    Open the mission map, find \textbf{Foundations 7 Capstone — Flight Computer Online}, and fly it for real ---
    everything you just saw, you now do yourself.
\end{frame}

\end{document}
